problem:  
Let \( W_{p,q} = \mathrm{SU}(3)/\mathrm{U}(1)_{p,q} \) be the Aloff–Wallach space defined by  
\[
\mathrm{U}(1)_{p,q} = \{ \operatorname{diag}(e^{ip\theta}, e^{iq\theta}, e^{-i(p+q)\theta}) : \theta \in \mathbb{R} \}.
\]  
Each \(W_{p,q}\) is a compact, simply connected 7‑manifold with \(H^4(W_{p,q},\mathbb{Z}) \cong \mathbb{Z}_{p^2 + pq + q^2}\).  
It is known that certain \(W_{p,q}\) admit homogeneous metrics with positive sectional curvature.  

Consider all Riemannian metrics \(g\) on \(W_{p,q}\) satisfying the **two-sided curvature bound**
\[
1 \le \sec_g \le C
\]
for some universal constant \(C>1\), and normalize so that \(\operatorname{Vol}(W_{p,q},g)=1\).  

Open question:  
Does there exist a universal constant \(\varepsilon = \varepsilon(C) > 0\) such that for *all coprime pairs* \((p,q)\) admitting such a metric, one has  
\[
\operatorname{diam}(W_{p,q}, g) \ge \varepsilon?
\]  
Equivalently, can one rule out the existence of a sequence of positively curved Aloff–Wallach 7‑manifolds \((W_{p_k,q_k}, g_k)\) with \(1 \le \sec_{g_k}\le C\) and unit volume that collapse in the Gromov–Hausdorff sense as \(|p_k|+|q_k| \to \infty\)?

Why is it a "good" problem:  
This corrected form poses a precise and open question at the intersection of **positive curvature geometry**, **collapsing phenomena**, and **topological complexity in homogeneous spaces**:

1. **Curvature versus topology:** Aloff–Wallach spaces form an infinite discrete family of simply connected 7‑manifolds distinguished by torsion in \(H^4\). The problem asks whether uniform two‑sided curvature bounds impose geometric rigidity across this infinitely many topological types, probing whether curvature can control or “ignore” the algebraic growth of \(H^4(W_{p,q})\).

2. **Natural link to bounded‑curvature collapse theory:** The inclusion of the bound \(1 \le \sec \le C\) situates the problem within the Cheeger–Gromov–Fukaya framework of collapsing with bounded curvature, making the equivalence with Gromov–Hausdorff convergence rigorous.

3. **Measure of geometric finiteness:** If such an \(\varepsilon(C)\) existed, it would indicate a kind of non‑collapsing finiteness for this class of highly structured positively curved manifolds. Failure of such a bound would suggest that curvature permits arbitrarily thin, yet bounded‑curvature, geometries distributed over increasingly complex topological families.

4. **Extreme simplicity and conceptual depth:** The statement is concise—essentially a question about a uniform diameter lower bound—but its resolution involves deep tools from Riemannian geometry, Lie group theory, and geometric topology. It functions as a test of how positive curvature interacts with infinite arithmetic families of manifolds, embodying the “narrow entrance, profound depth” ideal of a **good** research problem.